\begin{frame}{Caso integrador: Empresa Peruana ABC S.A.C.}
\casoeducativo
\small
\begin{columns}[T]
\column{0.5\textwidth}
\textbf{Activos}
\begin{tabular}{lr}
Efectivo & \soles{120,000.00}\\
Cuentas por cobrar & \soles{240,000.00}\\
Inventarios & \soles{180,000.00}\\
Activo corriente & \soles{540,000.00}\\
Activo no corriente & \soles{660,000.00}\\
\textbf{Total activo} & \soles{1,200,000.00}\\
\end{tabular}
\column{0.5\textwidth}
\textbf{Pasivos y patrimonio}
\begin{tabular}{lr}
Pasivo corriente & \soles{300,000.00}\\
Pasivo no corriente & \soles{270,000.00}\\
Pasivo total & \soles{570,000.00}\\
Patrimonio & \soles{630,000.00}\\
\textbf{Total} & \soles{1,200,000.00}\\
\end{tabular}
\end{columns}
\end{frame}

\begin{frame}{Estado de Resultados resumido}
\small
\begin{center}
\begin{tabular}{lr}
\toprule
Concepto & Importe \\
\midrule
Ventas netas & \soles{1,500,000.00}\\
Costo de ventas & \soles{975,000.00}\\
Utilidad bruta & \soles{525,000.00}\\
Gastos operativos & \soles{300,000.00}\\
Utilidad operativa & \soles{225,000.00}\\
Gastos financieros & \soles{45,000.00}\\
Utilidad antes de impuestos & \soles{180,000.00}\\
Impuesto a la renta & \soles{53,100.00}\\
\textbf{Utilidad neta} & \soles{126,900.00}\\
\bottomrule
\end{tabular}
\end{center}
\end{frame}

\begin{frame}{Ratios de liquidez del caso integrador}
\[
\text{Razón corriente}=\frac{540,000.00}{300,000.00}=1.80
\]
\[
\text{Prueba ácida}=\frac{540,000.00-180,000.00}{300,000.00}=1.20
\]
\textbf{Interpretación conjunta:} existe cobertura de corto plazo, aunque parte de la liquidez depende de inventarios.
\end{frame}

\begin{frame}{Ratios de solvencia del caso integrador}
\[
\text{Endeudamiento}=\frac{570,000.00}{1,200,000.00}\times 100=47.50\text{ \char37}
\]
\[
\text{Deuda-patrimonio}=\frac{570,000.00}{630,000.00}=0.90
\]
\textbf{Lectura:} casi la mitad de los activos se financia con terceros.
\end{frame}

\begin{frame}{Ratios de rentabilidad del caso integrador}
\small
\[
\text{Margen neto}=\frac{126,900.00}{1,500,000.00}\times 100=8.46\text{ \char37}
\]
\[
\text{ROA simple}=\frac{126,900.00}{1,200,000.00}\times 100=10.58\text{ \char37}
\quad
\text{ROE simple}=\frac{126,900.00}{630,000.00}\times 100=20.14\text{ \char37}
\]
Para mayor rigor, ROA y ROE deben usar saldos promedio cuando se cuente con saldos iniciales y finales.
\end{frame}

\begin{frame}{Diagnóstico integral}
\begin{tabular}{p{0.22\textwidth}p{0.68\textwidth}}
\toprule
Área & Diagnóstico \\
\midrule
Liquidez & Capacidad de corto plazo razonable, con dependencia parcial de inventarios.\\
Gestión & Requiere saldos promedio para evaluar cobranza e inventarios.\\
Solvencia & El 47.50 \char37{} de activos está financiado por terceros.\\
Rentabilidad & Genera utilidad, pero debe compararse con metas y sector.\\
\bottomrule
\end{tabular}
\end{frame}
